% !TeX root = ../main.tex
\section{问题重述}

\subsection{问题背景}

企业从供应商处采购两类零配件并装配为电子产品。零配件和装配过程均可能产生不合格品，企业还需承担检测、装配、市场调换、拆解和重新生产等成本。因而，企业首先需要判断供应商给出的次品率是否可信，再依据质量信息选择各生产环节的处理策略。

\subsection{问题一：抽样检测方案设计}

已知供应商声称一批零配件的次品率不超过标称值，要求在检测次数尽可能少的条件下设计抽样检测方案。当标称次品率为 $10\%$ 时，方案应分别给出在 $95\%$ 信度下认定实际次品率超过标称值并拒收该批、在 $90\%$ 信度下认定实际次品率不超过标称值并接收该批的具体结果，即确定抽样检测次数以及对应的接收与拒收判据。

\subsection{问题二：生产过程决策}

给定两类零配件和成品的次品率，以及采购、检测、装配、调换和拆解等成本参数，需要决定是否检测零配件、是否检测成品，以及检测出的不合格成品是否拆解。不检测的零配件直接进入装配；对退回的不合格品，重复相应的检测、拆解与装配过程。需要给出六种参数情形下的具体决策方案与决策依据。

\subsection{问题三：多工序生产过程决策}

把问题二扩展到多道工序、多个零配件的情形。给定图~1 所示的两道工序、8 个零配件、3 个半成品和 1 个成品的组装结构，以及表~2 给出的全部次品率、采购单价、检测成本、装配成本、拆解费用、市场售价和调换损失，需要确定各零配件是否检测、各半成品是否检测、不合格半成品是否拆解、成品是否检测以及不合格成品是否拆解，并给出具体决策方案、决策依据和相应指标。

\subsection{问题四：次品率带抽样不确定性的重决策}

在问题二与问题三中，零配件、半成品和成品的次品率均假设为通过抽样检测方法得到（例如问题一使用的抽样方案）。需要把表~1 与表~2 中的次品率视为带有抽样不确定性的估计值，在此基础上重新给出问题二各情形与问题三生产结构的决策方案、决策依据及相应指标。
